% Bagian: first-order-logic
% Bab: syntax-and-semantics
% Subbagian: first-order-languages

\documentclass[../../../include/open-logic-section]{subfiles}

\begin{document}

\olfileid[id]{fol}{syn}{fol}

\olsection{Bahasa Orde Pertama}


Ungkapan logika orde pertama dibangun dari kosakata dasar yang memuat
\emph{!!{variable}s}, \emph{!!{constant}s}, \emph{!!{predicate}s}, dan
kadang-kadang \emph{!!{function}s}. Dari semuanya itu, bersama penghubung
logis, kuantor, dan tanda baca seperti tanda kurung dan koma, dibentuk
\emph{suku} dan \emph{!!{formula}s}.

\begin{explain}
Secara informal, !!{predicate}s merupakan nama sifat dan relasi,
!!{constant}s merupakan nama objek individual, dan !!{function}s merupakan
nama pemetaan. Semua ini, kecuali !!{identity}~$\eq$, adalah
\emph{simbol nonlogis} dan secara bersama-sama membentuk suatu bahasa.
Setiap bahasa orde pertama~$\Lang L$ ditentukan oleh simbol-simbol
nonlogisnya. Dalam kasus yang paling umum, $\Lang L$ memuat tak hingga
banyaknya simbol dari setiap jenis.
\end{explain}

Dalam kasus umum, kita menggunakan simbol-simbol berikut dalam logika orde
pertama:

\begin{enumerate}
\item Simbol logis
\begin{enumerate}
\item Penghubung logis:
  \startycommalist
  \iftag{prvNot}{\ycomma $\lnot$ (negasi)}{}%
  \iftag{prvAnd}{\ycomma $\land$ (konjungsi)}{}%
  \iftag{prvOr}{\ycomma $\lor$ (disjungsi)}{}%
  \iftag{prvIf}{\ycomma $\lif$ (!!{conditional})}{}%
  \iftag{prvIff}{\ycomma $\liff$ (!!{biconditional})}{}%
  \iftag{prvAll}{\ycomma $\lforall$ (kuantor universal)}{}%
  \iftag{prvEx}{\ycomma $\lexists$ (kuantor eksistensial)}{}.
\tagitem{prvFalse}{Konstanta proposisional untuk !!{falsity}~$\lfalse$.}{}
\tagitem{prvTrue}{Konstanta proposisional untuk !!{truth}~$\ltrue$.}{}
\item !!{identity} dua-tempat~$\eq$.
\item Suatu himpunan !!{denumerable} dari !!{variable}s: $\Obj v_0$, $\Obj v_1$, $\Obj
  v_2$, \dots
\end{enumerate}
\item Simbol nonlogis yang membentuk \emph{bahasa standar} logika orde
  pertama
\begin{enumerate}
\item Suatu himpunan !!{denumerable} dari !!{predicate}s $n$-tempat untuk
  setiap $n>0$: $\Obj A^n_0$, $\Obj A^n_1$, $\Obj A^n_2$, \dots
\item Suatu himpunan !!{denumerable} dari !!{constant}s: $\Obj c_0$, $\Obj c_1$, $\Obj
  c_2$, \dots.
\item Suatu himpunan !!{denumerable} dari !!{function}s $n$-tempat untuk
  setiap $n>0$: $\Obj f^n_0$, $\Obj f^n_1$, $\Obj f^n_2$, \dots
\end{enumerate}
\item Tanda baca: (, ), dan koma.
\end{enumerate}

Sebagian besar definisi dan hasil kita akan dirumuskan untuk keseluruhan
bahasa standar logika orde pertama. Namun, bergantung pada penerapannya, kita juga
dapat membatasi bahasa itu hanya pada beberapa !!{predicate}s,
!!{constant}s, dan !!{function}s.

\begin{ex}
Bahasa aritmetika~$\Lang L_A$ memuat satu !!{predicate} dua-tempat~$<$,
satu !!{constant}~$\Obj 0$, satu !!{function} satu-tempat~$\prime$, serta dua
!!{function}s dua-tempat, yaitu~$+$ dan~$\times$.
\end{ex}

\begin{ex}
Bahasa teori himpunan~$\Lang L_Z$ hanya memuat satu !!{predicate}
dua-tempat~$\in$.
\end{ex}

\begin{ex}
Bahasa relasi urutan~$\Lang L_\le$ hanya memuat satu !!{predicate}
dua-tempat~$\le$.
\end{ex}

Sekali lagi, semua ini merupakan konvensi: secara resmi, semuanya hanyalah
alias. Misalnya, $<$, $\in$, dan $\le$ merupakan alias untuk $\Obj A^2_0$,
$\Obj 0$ untuk $\Obj c_0$, $\prime$ untuk $\Obj f^1_0$, $+$ untuk $\Obj f^2_0$,
dan $\times$ untuk $\Obj f^2_1$.

\iftag{defNot,defOr,defAnd,defIf,defIff,defTrue,defFalse,defEx,defAll}{%
Selain penghubung primitif dan
\iftag{notprvEx,notprvAll}{kuantor}{kuantor-kuantor} yang diperkenalkan di
atas, kita juga menggunakan simbol-simbol \emph{terdefinisi} berikut:
\startycommalist
  \iftag{defNot}{\ycomma $\lnot$ (negasi)}{}%
  \iftag{defAnd}{\ycomma $\land$ (konjungsi)}{}%
  \iftag{defOr}{\ycomma $\lor$ (disjungsi)}{}%
  \iftag{defIf}{\ycomma $\lif$ (!!{conditional})}{}%
  \iftag{defIff}{\ycomma $\liff$ (!!{biconditional})}{}%
  \iftag{defAll}{\ycomma $\lforall$ (kuantor universal)}{}%
  \iftag{defEx}{\ycomma $\lexists$ (kuantor eksistensial)}{}%
  \iftag{defFalse}{\ycomma !!{falsity}~$\lfalse$}{}%
  \iftag{defTrue}{\ycomma !!{truth}~$\ltrue$}{}.}{}

\begin{tagblock}{defNot,defOr,defAnd,defIf,defIff,defTrue,defFalse,defEx,defAll}
\begin{explain}
Suatu simbol terdefinisi secara resmi bukan bagian dari bahasa, tetapi
diperkenalkan sebagai singkatan informal: simbol itu memungkinkan kita
menyingkat formula yang akan menjadi cukup panjang jika kita hanya memakai
simbol primitif. Hal ini jelas merupakan suatu keuntungan. Namun, keuntungan
yang lebih besar ialah bahwa bukti menjadi lebih singkat. Jika suatu simbol
bersifat primitif, simbol itu harus ditangani secara tersendiri dalam bukti.
Oleh karena itu, semakin banyak simbol primitif, semakin panjang bukti kita.
\end{explain}
\end{tagblock}

% Simbol alternatif

\begin{intro}
Anda mungkin lebih mengenal istilah dan simbol yang berbeda dari yang kita
gunakan di atas. Teks-teks logika (dan para pengajar) lazim menggunakan
$\sim$, $\neg$, atau~!\ untuk ``negasi'', serta $\wedge$, $\cdot$, atau $\&$
untuk ``konjungsi''. Simbol yang lazim digunakan untuk ``implikasi'' ialah
$\rightarrow$, $\Rightarrow$, dan $\supset$.
\iftag{prvIff,defIff}{Simbol untuk ``biimplikasi'', ``implikasi dua arah'',
  atau ``ekuivalensi (material)'' ialah $\leftrightarrow$,
  $\Leftrightarrow$, dan $\equiv$.}{}
\iftag{prvFalse,defFalse}{Simbol $\lfalse$ disebut dengan berbagai nama,
  antara lain ``kepalsuan'', ``falsum'', ``absurditas'', atau ``bawah''.}{}
\iftag{prvTrue,defTrue}{Simbol $\ltrue$ disebut dengan berbagai nama,
  antara lain ``kebenaran'', ``verum'', atau ``atas''.}{}

Secara konvensional, huruf kecil (misalnya, $a$, $b$, $c$) dari awal alfabet
Latin digunakan untuk !!{constant}s (kadang-kadang disebut nama), sedangkan
huruf kecil dari akhir alfabet (misalnya, $x$, $y$, $z$) digunakan untuk
!!{variable}s. Kuantor digabungkan dengan !!{variable}s, misalnya $x$;
variasi notasinya mencakup $\forall x$, $(\forall x)$, $(x)$, $\Pi x$,
$\bigwedge_x$ untuk kuantor universal dan $\exists x$, $(\exists x)$,
$(Ex)$, $\Sigma x$, $\bigvee_x$ untuk kuantor eksistensial.
\end{intro}

\begin{explain}
Kita dapat memperlakukan semua operator proposisional dan kedua kuantor
sebagai simbol primitif bahasa. Sebagai gantinya, kita dapat memilih
persediaan simbol primitif yang lebih kecil dan memperlakukan !!{operator}s
lain sebagai simbol terdefinisi. Himpunan operator Boolean yang ``lengkap
secara fungsional-kebenaran'' mencakup $\{ \lnot, \lor \}$, $\{ \lnot,
\land \}$, dan $\{ \lnot, \lif\}$---masing-masing dapat digabungkan dengan
salah satu kuantor untuk menghasilkan bahasa orde pertama yang lengkap
secara ekspresif.

Anda mungkin mengenal dua !!{operator}s lain: garis Sheffer~$|$ (dinamai
menurut Henry Sheffer) dan panah Peirce~$\downarrow$, yang juga dikenal
sebagai belati Quine. Dengan pembacaan lazimnya masing-masing sebagai
``nand'' dan ``nor'', kedua operator ini secara sendiri-sendiri lengkap
secara fungsional-kebenaran.
\end{explain}

\end{document}
